\section{Formal Verification}
\label{sec:verification}
%% ===================================================================

All critical protocols in the trust kernel have corresponding Lean~4 mechanized
proofs. The proof modules are organized as follows:

\begin{table}[h]
\centering
\caption{Lean~4 proof modules for trust kernel properties.}
\label{tab:proofs}
\begin{tabular}{lp{7cm}}
\toprule
\textbf{Module} & \textbf{Property} \\
\midrule
\texttt{Consensus/BFT} & Safety and liveness under $f < n/3$ Byzantine
    validators \\
\texttt{Crypto/FROST} & Unforgeability of FROST threshold Schnorr
    signatures \\
\texttt{Crypto/Threshold/CGGMP21} & Threshold ECDSA composition security
    and identifiable abort correctness \\
\texttt{Trust/Authority} & Trust delegation chain integrity: no capability
    escalation, revocation finality \\
\bottomrule
\end{tabular}
\end{table}

\begin{remark}
The Lean~4 proofs verify the \emph{protocol logic}, not the implementation.
Implementation correctness is established through the type-safe Go and Rust
codebases (\texttt{luxfi/frost}, \texttt{luxfi/threshold}, \texttt{luxfi/lattice})
with property-based testing against the Lean~4 reference specifications.
\end{remark}

The proof of trust delegation integrity (Theorem~\ref{thm:signing-auth}) is
mechanized in \texttt{Trust/Authority.lean} as follows: the proof establishes
that the authorization graph is acyclic, that revocation propagates to all
downstream delegations, and that no delegation chain can exceed the
capabilities of its root authority.

%% ===================================================================
